计划的180个方案与随机状态组合中，实际启动180次后处理；未取得共享LP最优解的组合记为未启动，不填零耗时。20次共享LP准备中，20次取得最优状态，平均耗时为5.12秒，样本标准差为5.91秒。LP-G和LP-N若作为独立端到端方法使用，还须各自计入相应输入的LP准备成本；UNI无需该阶段。

LP-G在计划的60组后处理设置中取得55次可行方案。相同可行运行上，局部优化前后平均遗漏分别为90.2和64.1，其中55次严格改善。该变化直接来自后续局部搜索，不能归因于初次LP分配。

LP-N在计划的60组后处理设置中取得56次可行方案。相同可行运行上，局部优化前后平均遗漏分别为93.9和65.6，其中56次严格改善。该变化直接来自后续局部搜索，不能归因于初次LP分配。

UNI在计划的60组后处理设置中取得56次可行方案。相同可行运行上，局部优化前后平均遗漏分别为101.9和63.1，其中56次严格改善。该变化直接来自后续局部搜索，不能归因于初次LP分配。

在LP-G与LP-N共同可行的20个独立输入上，LP-G取得5优、4平、11劣，Holm校正后$p=0.4304$。本次条件配对检验未在0.05水平下给出分量引导改善目标质量的显著证据；这并不证明两种方案等效。
